\section{Blast vs SuperFaB}
%==============================================================

Both \textbf{Blast.jl} \cite{chiarenza2024blast} and \textbf{SuperFaB}
\cite{grasshorn2022superfab} are modern Julia codes designed to compute
statistics related to the angular (or spherical) power spectrum
$C_\ell$. Despite this shared goal, the two codes address fundamentally
different scientific problems:

\begin{itemize}
    \item \textbf{Blast.jl} is a \emph{theory prediction code}.
    It computes the projected 2D angular power spectrum
    $C_\ell^{AB}$ between tomographic bins of galaxy surveys without
    invoking the Limber approximation, allowing sub-percent accuracy
    at low multipoles.

    \item \textbf{SuperFaB} is a \emph{data estimator}. It measures
    the full 3D Spherical Fourier-Bessel (SFB) power spectrum
    $C_\ell(k, k')$ directly from a galaxy catalogue, retaining all
    radial and angular information without collapsing it into redshift
    bins.
\end{itemize}

The two approaches are complementary: SuperFaB provides a richer
estimator suited for all-sky spectroscopic surveys, while Blast.jl
provides a fast and accurate theory model especially relevant for
tomographic (photometric) analyses such as 3$\times$2\,pt statistics.

%==============================================================
\subsection{The Angular Power Spectrum: Two Perspectives}
%==============================================================

\subsection{Blast.jl --- Beyond-Limber Angular Power Spectra}

\subsubsection{From 3D to 2D: the projection integral}

Any 2D field observed on the sky is a projection of a 3D field
$A(\mathbf{r}, z)$ through a window function $W_i(z)$:
\begin{equation}
    \tilde{A}(\hat{n}) = \int dz\, W_i(z)\, A(\chi(z)\hat{n};\, z),
\end{equation}
where $\chi(z)$ is the comoving distance. Expanding in spherical
harmonics via the Rayleigh formula $e^{i\mathbf{k}\cdot\hat{n}r} =
4\pi\sum_{\ell m}i^\ell j_\ell(kr)Y^*_{\ell m}(\hat{k})Y_{\ell m}(\hat{n})$,
one obtains the angular power spectrum between probes $A$ and $B$ in
tomographic bins $i$ and $j$:
\begin{equation}
    \boxed{
    C_{ij}^{AB}(\ell) = \mathcal{N}(\ell)
    \int_0^\infty d\chi_1\, W_i^A(\chi_1)
    \int_0^\infty d\chi_2\, W_j^B(\chi_2)
    \int_0^\infty dk\, k^2\,
    \frac{P^{AB}(k,\chi_1,\chi_2)\,j_\ell(k\chi_1)\,j_\ell(k\chi_2)}
         {(k\chi_1)^\alpha (k\chi_2)^\beta}
    }
    \label{eq:blast_cij}
\end{equation}
where $j_\ell(x)$ are the spherical Bessel functions,
$P^{AB}(k,\chi_1,\chi_2)$ is the (possibly unequal-time) 3D matter
power spectrum, and the exponents $(\alpha,\beta)$ depend on the probe
combination:
\begin{equation}
    (\alpha, \beta) =
    \begin{cases}
        (0,0) & \text{galaxy clustering (gg)}\\
        (2,2) & \text{weak lensing shear (ss)}\\
        (0,2) & \text{galaxy--galaxy lensing (gs)}
    \end{cases}
\end{equation}
The prefactor is $\mathcal{N}(\ell) = 2/\pi$ for clustering and
$\mathcal{N}(\ell) = 2(ell+2)!/[\pi(\ell-2)!]$ for lensing.

The probe-specific kernels are
\begin{align}
    W_i^g(\chi) &= \frac{H(z)}{c}\,n_i(z)\,b_g(z), \\
    W_i^s(\chi) &= \frac{3H_0^2\Omega_m}{2c^2\,a(\chi)}\,\chi
                   \int_z^\infty dz'\,n_i(z')
                   \frac{\chi(z')-\chi}{\chi(z')}.
\end{align}

\subsubsection{The Limber approximation and its breakdown}

The standard Limber approximation replaces the spherical Bessel
function by a Dirac delta,
\begin{equation}
    j_\ell(x) \approx \sqrt{\frac{\pi}{2\ell+1}}\,
    \delta^D\!\left(\ell+\tfrac{1}{2} - x\right),
\end{equation}
collapsing the double $(\chi_1,\chi_2)$ integral to a single 1D
integral:
\begin{equation}
    C_{ij}^{AB}(\ell) \approx
    \int_0^\infty \frac{d\chi}{\chi^2}\,
    K_i^A(\chi)\,K_j^B(\chi)\,
    P^{AB}\!\left(k_\ell = \frac{\ell+1/2}{\chi},\,z\right).
    \label{eq:limber}
\end{equation}
While computationally cheap, Eq.~\eqref{eq:limber} breaks down for
$\ell \lesssim 100$, narrow kernels (spectroscopic bins), and
cross-bin correlations, introducing errors of order $\mathcal{O}(\ell^{-2})$.

\subsubsection{The Blast.jl Chebyshev algorithm}

The key idea of Blast.jl is to decompose the power spectrum on a
Chebyshev basis:
\begin{equation}
    P(k,\chi_1,\chi_2;\,\theta) \approx
    \sum_{n=0}^{n_{\max}} c_n(\chi_1,\chi_2;\,\theta)\,T_n(k),
    \label{eq:cheb_pk}
\end{equation}
where $T_n(k)$ are Chebyshev polynomials of the first kind and $c_n$
are coefficients that carry all cosmological parameter dependence.
This separates the integral in Eq.~\eqref{eq:blast_cij} into a
cosmology-independent (precomputable) part and a fast update:
\begin{equation}
    w_\ell^{AB}(\chi_1,\chi_2;\,\theta) =
    \sum_{n=0}^{n_{\max}} c_n(\chi_1,\chi_2;\,\theta)\,
    \tilde{T}_{n;\,\ell}^{AB}(\chi_1,\chi_2),
    \label{eq:w_decomp}
\end{equation}
where the cosmology-\emph{independent} precomputed integrals are:
\begin{equation}
    \tilde{T}_{n;\,\ell}^{AB}(\chi_1,\chi_2) \equiv
    \int_{k_{\min}}^{k_{\max}} dk\,f^{AB}(k)\,T_n(k)\,
    j_\ell(k\chi_1)\,j_\ell(k\chi_2),
\end{equation}
with $f^{gg}(k)=k^2$, $f^{ss}(k)=k^{-2}$, $f^{gs}(k)=1$.

Introducing the coordinate change $R \equiv \chi_2/\chi_1$, the
outer double integral becomes
\begin{equation}
    C_{ij}^{AB}(\ell) =
    \int_0^\infty d\chi
    \int_0^1 dR\;\chi\,
    \Bigl[K_i^A(\chi)\,K_j^B(R\chi)
         + K_j^B(\chi)\,K_i^A(R\chi)\Bigr]\,
    w_\ell^{AB}(\chi,\,R\chi),
\end{equation}
which is evaluated with Simpson quadrature in $\chi$ and
Clenshaw-Curtis quadrature in $R$ (48 Chebyshev points in $(0,1]$),
where the sharp feature at $R=1$ ($\chi_1 \approx \chi_2$) is well
resolved by the Chebyshev node distribution.

\subsubsection{Blast.jl implementation summary}

The code stores three precomputed tensor arrays of $\tilde{T}$ (one
per probe combination, with $\beta=\{2,0,-2\}$) loaded at module
initialisation from \texttt{artifact} files. At inference time:
\begin{enumerate}
    \item Evaluate Chebyshev coefficients $c_n$ via FFT (DCT-I) of
          $P(k_j,\chi_1,\chi_2)$ on the Chebyshev grid.
    \item Contract $c_n$ with $\tilde{T}_{n;\ell}$ via \texttt{@tullio}
          (Einstein notation) to obtain $w_\ell^{AB}(\chi,R)$.
    \item Build the kernel product $K^{[2]}_{ij}(\chi,R)$ via a
          \texttt{combine\_kernels} call.
    \item Integrate with Simpson ($\chi$) and Clenshaw-Curtis ($R$)
          weights to obtain $C_{ij}^{AB}(\ell)$.
\end{enumerate}
On a 32-core node, the fiducial LSST-Y10 data vector (10 clustering
bins + 5 lensing bins, $\ell\in[2,200]$) is computed in $\approx 3.6$\,ms,
roughly 10--15$\times$ faster than the previous best code (FKEM) in the
N5K challenge \cite{leonard2023n5k}.

%------------------------------------------------------------
\subsection{SuperFaB --- Spherical Fourier-Bessel Decomposition}
%------------------------------------------------------------

\subsubsection{The SFB transform}

SuperFaB works directly in 3D. An overdensity field
$\delta(\mathbf{r})$ is expanded in the SFB basis:
\begin{equation}
    \delta(\mathbf{r}) =
    \sum_{n\ell m}\bigl[g_{n\ell}(r)\,Y_{\ell m}(\hat{r})\bigr]
    \delta_{n\ell m},
    \qquad
    \delta_{n\ell m} =
    \int d^3r\,\bigl[g_{n\ell}(r)\,Y_{\ell m}^*(\hat{r})\bigr]
    \delta(\mathbf{r}),
\end{equation}
with the integration over the volume $r_{\min}\le r\le r_{\max}$.
The radial basis functions $g_{n\ell}(r)$ are linear combinations of
spherical Bessel functions of the first and second kind,
\begin{equation}
    g_{n\ell}(r) = c_{n\ell}\,j_\ell(k_{n\ell}\,r)
                 + d_{n\ell}\,y_\ell(k_{n\ell}\,r),
\end{equation}
satisfying \emph{potential boundary conditions} at $r_{\min}$ and
$r_{\max}$ (Fisher et al.~1995, Samushia 2019):
\begin{equation}
    \ell\, g_{n\ell}(k_{n\ell}r_{\min})
    = k_{n\ell}r_{\min}\,g'_{n\ell}(k_{n\ell}r_{\min}),
    \quad
    -({\ell+1})\,g_{n\ell}(k_{n\ell}r_{\max})
    = k_{n\ell}r_{\max}\,g'_{n\ell}(k_{n\ell}r_{\max}),
\end{equation}
and the orthonormality condition:
\begin{equation}
    \int_{r_{\min}}^{r_{\max}} dr\,r^2\,
    g_{n\ell}(r)\,g_{n'\ell}(r) = \delta_{nn'}^K.
\end{equation}

\subsubsection{The SFB power spectrum}

For a survey with radial selection function $\phi(r)$ and linear
growth $D(r)$, the observable SFB window becomes:
\begin{equation}
    \mathcal{W}_\ell(k, q) =
    \frac{2qk}{\pi}
    \int dr\,r^2\,\phi(r)\,D(r)\,b(r,q)\,
    j_\ell(kr)\,j_\ell(qr).
    \label{eq:sfb_window}
\end{equation}
The SFB power spectrum is then:
\begin{equation}
    \boxed{
    C_\ell(k, k') = \int dq\,
    \mathcal{W}_\ell(k, q)\,\mathcal{W}^*_\ell(k', q)\,P(q).
    }
    \label{eq:sfb_cll}
\end{equation}
This is a \emph{matrix} in $(k,k')$: diagonal entries
$C_\ell(k,k)$ carry the auto-power at scale $k$, while off-diagonal
entries encode mode coupling induced by redshift evolution of the
growth factor and bias.

In the Limber approximation $[j_\ell(kr)\approx
\sqrt{\pi/(2\ell+1)}\,\delta^D(k-\ell+1/2)/r]\,r^{-1}$, the SFB
window simplifies to
\begin{equation}
    \mathcal{W}_\ell^{\rm Limb}(k,q) =
    \delta^D(q-k)\,\phi\!\left(\tfrac{\ell+1/2}{k}\right)
    D\!\left(\tfrac{\ell+1/2}{k}\right)
    b\!\left(\tfrac{\ell+1/2}{k},k\right),
\end{equation}
so $C_\ell(k,k')\to\delta^D(k-k')\,C_\ell(k)$, i.e.\ the spectrum
becomes diagonal and the full 3D information collapses to a 1D
function of $k$ at the characteristic radius $r=(\ell+1/2)/k$.

\subsubsection{Estimation from data}

SuperFaB does not compute a theory model but \emph{measures}
Eq.~\eqref{eq:sfb_cll} from galaxies. The 3DEX approach
(Leistedt et al.\ 2012) first performs the radial transform:
\begin{equation}
    \delta_{n\ell}^{\rm obs}(\hat{r}) =
    \frac{1}{\bar{n}}\sum_p
    \delta^D(\hat{r}-\hat{r}_p)\,g_{n\ell}(r_p)
    - \mathcal{W}_{n\ell}(\hat{r}),
\end{equation}
then the angular spherical harmonic transform via HEALPix.
The pseudo-SFB power spectrum estimator is:
\begin{equation}
    \hat{C}_\ell^{\rm obs}(n,n') =
    \frac{1}{2\ell+1}\sum_m
    \delta_{n\ell m}^{\rm obs}
    \bigl(\delta_{n'\ell m}^{\rm obs}\bigr)^*,
\end{equation}
which is then corrected for the window mixing matrix
$\mathcal{M}_{\ell nn'}^{LNN'}$, shot noise, and pixel window.

%==============================================================
\subsection{Key Conceptual Differences}
%==============================================================

\subsection{Observables and dimensionality}

\begin{center}
\renewcommand{\arraystretch}{1.5}
\begin{tabular}{lcc}
\toprule
& \textbf{Blast.jl} & \textbf{SuperFaB} \\
\midrule
Output observable & $C_\ell^{AB}$ (scalar per $\ell$, bin pair) &
    $C_\ell(k,k')$ (matrix in $k$) \\
Info retained & Projected 2D & Full 3D radial + angular \\
Survey type & Photometric (tomographic) & Spectroscopic / all-sky \\
Role & Theory predictor & Data estimator \\
Limber regime & Goes beyond & Recovers Limber on diagonal \\
\bottomrule
\end{tabular}
\end{center}

\subsection{Relationship between $C_\ell^{AB}$ and $C_\ell(k,k')$}

The Blast.jl projected angular power spectrum can be seen as a
\emph{weighted integral of the SFB power spectrum}. Defining
$\hat{W}_i^A(k) \equiv \int d\chi\,W_i^A(\chi)\,j_\ell(k\chi)$,
one has schematically:
\begin{equation}
    C_{ij}^{AB}(\ell) \sim
    \int dk\, dk'\,
    \hat{W}_i^A(k)\,\hat{W}_j^B(k')\,
    C_\ell(k, k').
    \label{eq:connection}
\end{equation}
In the Limber limit $C_\ell(k,k')\to\delta^D(k-k')C_\ell(k)$, this
becomes Eq.~\eqref{eq:limber}. Beyond Limber, both off-diagonal terms
in $C_\ell(k,k')$ and the exact spherical Bessel integrals in
$\hat{W}_i^A(k)$ must be retained --- this is precisely what
Blast.jl does via its Chebyshev decomposition in $k$.

Alternatively, the diagonal slice of the SFB spectrum contains a
theory prediction:
\begin{equation}
    C_\ell(k,k) =
    \int dq\,|\mathcal{W}_\ell(k,q)|^2\,P(q)
    \xrightarrow{\rm Limber}
    \phi^2\!\left(\frac{\ell+1/2}{k}\right)
    D^2\!\left(\frac{\ell+1/2}{k}\right)
    b^2\!\left(\frac{\ell+1/2}{k},k\right)
    P(k),
\end{equation}
which is entirely analogous to the Limber approximation of $C_{ij}^{AB}$
for a single broad bin with $W_i(\chi)=\phi(\chi)D(\chi)b(\chi)$.

\subsection{Treatment of the spherical Bessel functions}

Both codes must handle products of spherical Bessel functions
$j_\ell(k\chi_1)j_\ell(k\chi_2)$, but they do so in opposite
strategies:

\begin{center}
\renewcommand{\arraystretch}{1.5}
\begin{tabular}{lcc}
\toprule
& \textbf{Blast.jl} & \textbf{SuperFaB} \\
\midrule
Bessel integration & Pre-computed $\tilde{T}_{n\ell}^{AB}$ tensors &
    Discrete $k_{n\ell}$ modes from BCs \\
Basis & Chebyshev $T_n(k)$ & Eigenmodes $g_{n\ell}(r)$ of $\nabla^2$ \\
$k$-grid & Clenshaw-Curtis ($2^{15}+1$ pts) & Discrete $\{k_{n\ell}\}$ \\
Oscillations handled by & High-$N$ CC quadrature & Orthogonality of $g_{n\ell}$ \\
\bottomrule
\end{tabular}
\end{center}

\subsection{Coordinate system for the outer integrals}

Both codes use a change of coordinates to improve numerical accuracy:

\begin{itemize}
    \item \textbf{Blast.jl}: $(\chi_1,\chi_2)\to(\chi, R=\chi_2/\chi_1)$,
    concentrating quadrature points near $R=1$ (same epoch,
    diagonal of $P(k,\chi_1,\chi_2)$).
    \item \textbf{SuperFaB}: $(r,\hat{r})\to(n\ell, \hat{r})$,
    separating the radial transform (performed first) from the angular
    HEALPix transform; the boundary at $r_{\min}$ eliminates the
    numerically unstable empty volume.
\end{itemize}

\subsection{Boundary conditions and mode discreteness}

A distinctive feature of SuperFaB is the imposition of potential
boundary conditions at \emph{both} $r_{\min}$ and $r_{\max}$.
This discretises the allowed wavenumbers $\{k_{n\ell}\}$ and
orthogonalises the basis functions $g_{n\ell}$, making the window
deconvolution ($\mathcal{M}^{-1}$) numerically stable.
Blast.jl has no analogous discretisation: it works with a continuous
$k$-range $[k_{\min}, k_{\max}]$ sampled on Chebyshev points for the
precomputed $\tilde{T}$ integrals.

%==============================================================
\subsection{Comparison of the $C_\ell$ Result}
%==============================================================

\subsection{Limber approximation as a common benchmark}

In the Limber limit, both formalisms converge:

\begin{align}
    C_{ij}^{AB,\,\rm Limber}(\ell) &=
    \int \frac{d\chi}{\chi^2}\,
    K_i^A(\chi)\,K_j^B(\chi)\,
    P\!\left(\frac{\ell+1/2}{\chi},z(\chi)\right),
    \label{eq:Clim_blast}\\[6pt]
    C_\ell^{\rm Limber}(k) &=
    \phi^2\!\left(\frac{\ell+1/2}{k}\right)
    D^2\!\left(\frac{\ell+1/2}{k}\right)
    b^2\!\left(\frac{\ell+1/2}{k},k\right)
    P(k).
    \label{eq:Clim_sfb}
\end{align}

These are related by $C_{ij}^{AB,\,\rm Limber}(\ell) =
\int dk\,\hat{W}_i^A(k)^2\,C_\ell^{\rm Limber}(k)$
for the auto-spectrum of a single bin.

\subsection{Beyond-Limber corrections}

Going beyond Limber reveals different structures in the two formalisms:

\textbf{Blast.jl (beyond-Limber $C_\ell^{AB}$):}
The full 3D integral in Eq.~\eqref{eq:blast_cij} captures:
\begin{itemize}
    \item Non-zero off-diagonal $(\chi_1\ne\chi_2)$ contributions,
    which are suppressed only as $\mathcal{O}(\ell^{-2})$ in the Limber
    limit.
    \item Exact unequal-time power spectrum
    $P(k,\chi_1\ne\chi_2)$, important for cross-bin correlations and
    at low $\ell$.
    \item The Chebyshev approach recovers sub-percent accuracy already
    at $\ell = 2$ with $n_{\max}=120$ basis functions.
\end{itemize}

\textbf{SuperFaB (full SFB $C_\ell(k,k')$):}
The off-diagonal $C_\ell(k\ne k')$ terms encode redshift evolution:
\begin{equation}
    C_\ell(k,k') \xrightarrow{{\rm evol.}}
    \phi^2\!\left(\frac{\ell+1/2}{k}\right)
    \phi^2\!\left(\frac{\ell+1/2}{k'}\right)
    D\!\left(\frac{\ell+1/2}{k}\right)
    D\!\left(\frac{\ell+1/2}{k'}\right)
    \int dq\,P(q)\,\delta^D(q-k)\,\delta^D(q-k').
\end{equation}
Redshift evolution of $D(r)$ and $b(r)$ tilts the relation between
the measured $k$-mode and the effective redshift, populating the
off-diagonal elements of $C_\ell(k,k')$.

\subsection{Accuracy and practical performance}

\begin{center}
\renewcommand{\arraystretch}{1.5}
\begin{tabular}{lcc}
\toprule
\textbf{Metric} & \textbf{Blast.jl} & \textbf{SuperFaB} \\
\midrule
N5K $\Delta\chi^2$ ($\ell\in[2,200]$) & $0.1054$ (pass, req.\ $<0.2$) &
    N/A (not an angular $C_\ell$ code) \\
Runtime (fiducial, 64 threads) & $5.12 \pm 0.01$\,ms &
    $\sim10$\,CPU-min/sim (Roman) \\
Pre-computation & $\tilde{T}$ once on laptop ($\sim$1-2 h) &
    BCs + $g_{n\ell}$ once per survey \\
Accuracy control & $n_{\max}$, $N_{\rm CC}$ & $\Delta\ell$, $\Delta n$ \\
$\chi^2_\nu$ on simulations & --- & $1.033$ (Roman, 213 modes) \\
Main limitation & $P(k,\chi_1,\chi_2)$ separability assumed
    for NL part & $\chi^2_\nu$ degraded for poorly-deconvolved modes \\
\bottomrule
\end{tabular}
\end{center}

\subsection{Formal connection via the SFB window matrix}

The SFB window matrix (Eq.~\eqref{eq:sfb_cll}) and the Blast.jl
kernel product are mathematically linked. Defining the SFB projection
of the galaxy kernel as
\begin{equation}
    \hat{W}_i^g(k;\ell) \equiv
    \frac{2}{\pi}k^2
    \int d\chi\,W_i^g(\chi)\,j_\ell(k\chi),
\end{equation}
the Blast.jl clustering power spectrum in Eq.~\eqref{eq:blast_cij} can
be rewritten as:
\begin{equation}
    C_{ij}^{gg}(\ell) =
    \int \frac{dk}{k^2}\,
    \hat{W}_i^g(k;\ell)\,
    \hat{W}_j^g(k;\ell)\,
    P(k)
    \quad + \quad \text{unequal-time terms},
\end{equation}
which is the projection of the SFB window $\mathcal{W}_\ell(k,q)$
(Eq.~\eqref{eq:sfb_window}) integrated over $q$.

This shows that the SuperFaB SFB window $\mathcal{W}_\ell$ and the
Blast.jl spherical Bessel integrals $\hat{W}_i^g$ are two faces of
the same mathematical object, differentiated only by whether the full
radial basis $g_{n\ell}$ (discrete, SuperFaB) or the continuous
spherical Bessel integral (Blast.jl) is used.

%==============================================================
\subsection{When to Use Each Code}
%==============================================================

\begin{center}
\renewcommand{\arraystretch}{1.4}
\begin{tabular}{p{5cm}p{4cm}p{4cm}}
\toprule
\textbf{Science goal} & \textbf{Recommended} & \textbf{Why} \\
\midrule
Parameter inference with tomographic
photometric data (3$\times$2\,pt) &
    Blast.jl & Fast theory model, MCMC-friendly \\
Primordial non-Gaussianity ($f_{\rm NL}$)
via scale-dep.\ bias in spectroscopic survey &
    SuperFaB & Retains large radial modes \\
Cross-correlation of photo-$z$ bins
at low $\ell$ &
    Blast.jl & Beyond-Limber accuracy at $\ell<100$ \\
RSD modelling, all wide-angle effects &
    SuperFaB & Line-of-sight per galaxy, Kaiser built in \\
Nonlinear scales ($k>0.3\,h/\mathrm{Mpc}$) &
    Blast.jl + Limber for NL part &
    NL contribution decays as $e^{-(D_1-D_2)^2}$ \\
\bottomrule
\end{tabular}
\end{center}

%==============================================================
\subsection{Connection to the Internship Code}
%==============================================================

The Julia source files provided (\texttt{blast\_tutorials.jl},
\texttt{integrals.jl}, \texttt{projected\_matter.jl},
\texttt{chebcoefs.jl}, \texttt{background.jl}) constitute the
prototype of Blast.jl. The key data-flow is:

\begin{enumerate}
    \item \texttt{background.jl}: computes $H(z)$ and $\chi(z)$
    by numerical quadrature of $1/E(z)$, populating
    \texttt{BackgroundQuantities}.
    \item \texttt{background.jl}: computes $W_i^g(\chi)$ and
    $W_i^s(\chi)$ using \texttt{compute\_kernel!} for
    \texttt{GalaxyKernel} and \texttt{ShearKernel} structs.
    \item \texttt{chebcoefs.jl}: computes fast Chebyshev
    coefficients $c_n$ via a pre-planned DCT-I FFT,
    implementing Eq.~\eqref{eq:cheb_pk}.
    \item \texttt{projected\_matter.jl}: computes $\tilde{T}_{n\ell}$
    in \texttt{compute\_T\~{}} using Clenshaw-Curtis quadrature and
    \texttt{LoopVectorization.@tturbo} inner loops, then contracts
    with Chebyshev coefficients in \texttt{w\_ell\_tullio}.
    \item \texttt{integrals.jl}: combines probe kernels in
    \texttt{combine\_kernels} (via \texttt{@tullio} Einstein
    contraction) and performs the $(\chi, R)$ outer integration in
    \texttt{compute\_C\ell}, recovering $C_{ij}^{AB}(\ell)$.
\end{enumerate}

The goal of the internship project is to extend this framework toward
the SFB decomposition of SuperFaB, i.e.\ to compute
$C_\ell(k,k')$ beyond the Limber approximation using the Blast.jl
Chebyshev machinery, providing a \emph{theory model} for the SFB
estimator that SuperFaB measures from data.

%==============================================================
\subsection{Conclusions}
%==============================================================

Blast.jl and SuperFaB attack the same family of integrals from
orthogonal directions:

\begin{itemize}
    \item \textbf{Blast.jl} precomputes the $k$-integrals of two
    Bessel functions against Chebyshev polynomials, then evaluates
    the outer $(\chi,R)$ integrals with quadrature at inference time.
    The result is a projected angular power spectrum $C_\ell^{AB}$
    valid beyond the Limber approximation.

    \item \textbf{SuperFaB} defines an orthonormal radial basis
    $g_{n\ell}$ from the Laplacian eigenfunctions with potential
    boundary conditions, measures the SFB coefficients from galaxy
    positions, and estimates the full 3D power spectrum matrix
    $C_\ell(k,k')$.
\end{itemize}

In the Limber limit, the diagonal of $C_\ell(k,k)$ from SuperFaB
and the $C_\ell^{AB}$ from Blast.jl converge to the same expression
(Eqs.~\eqref{eq:Clim_blast}--\eqref{eq:Clim_sfb}).
Beyond Limber, the Blast.jl approach provides exact projected spectra
for tomographic analysis, while SuperFaB retains the full 3D
structure needed for spectroscopic surveys sensitive to large radial
modes, RSD, and primordial non-Gaussianity.

The internship project bridges these two worlds: using the Chebyshev
precomputation strategy of Blast.jl to build a fast theory predictor
for the SFB power spectrum $C_\ell(k,k')$, enabling beyond-Limber
parameter inference in the SFB framework.